% GENERATED LOCAL-BUILD FLAT PART — DO NOT MERGE

% ---- BEGIN TIR/monograph/v12/chapters/ch01_first_distinction_relational_kernel.tex ----
% !TEX root = ../../tir_monograph_v12.tex
\chapter{First Distinction and the Relational Kernel}
\label{ch:v12_first_distinction}

\section{Dependency before temporal order}

The monograph begins before any object is admitted. Write
\begin{equation}
\boxed{X\prec_{\rm dep}Y}
\end{equation}
when the admitted structure represented by $Y$ uses $X$ as an upstream parent in the derivation graph. Temporal ordering receives its own dynamical treatment in the companion time branch; the primitive relation used here is therefore a dependency order.

TIR defines existential admissibility relationally. A relational presentation is
\begin{equation}
\mathfrak R=(X,\mathcal D),
\end{equation}
with support set $X$ and realized distinction/relation set $\mathcal D$. The pre-object zero layer is
\begin{equation}
\boxed{
\mathfrak Z_{\rm rel}=(\varnothing,\varnothing).
}
\label{eq:v12_relational_zero}
\end{equation}
It is not a point and not a singleton. It is zero realized relational structure. Therefore the support set contains no element.

The logical status is exact but scoped. Once TIR defines existence through realized relation or distinction, the absence of an existent object in the empty relational presentation follows by definition. The relational-existence rule itself is the declared TIR ontological entry condition.

Leaving relational zero requires nonempty support. The least nonempty support candidate is the point
\begin{equation}
\boxed{
\mathcal P=\{p\},
\qquad
|\mathcal P|=1.
}
\label{eq:v12_minimal_point_support}
\end{equation}
Thus
\begin{equation}
\boxed{
\mathfrak Z_{\rm rel}\neq\mathcal P.
}
\end{equation}
The unresolved point $(\mathcal P,\varnothing)$ supplies support but not yet a realized relational distinction.
\vTwelveStatus{B}{--}{PASS}

\section{Primitive kernel}

The A0 relational-existence root and the eight A1--A8 TIR axioms are grouped by role rather than by historical order:
\begin{description}
\item[A0 -- relational existence / zero:] defines the empty pre-object relational presentation $\mathfrak Z_{\rm rel}$ and separates zero from the later point support.
\item[A1 -- point minimality:] supplies the minimal non-empty support candidate $\mathcal P$.
\item[A2 -- quantum point:] supplies a normalized Hilbert-space representation for the minimal carrier.
\item[A3 -- information primacy:] supplies distinguishable informational relations and their normalized information carrier.
\item[A4 -- spherical geometric efficiency:] selects the sphere for the declared isotropic boundary/volume functional.
\item[A5 -- arithmetic measures geometry:] reads geometric closure through discrete invariants such as winding or multiplicity.
\item[A6 -- natural numbers from complex phase closure:] assigns natural closure indices to admitted complex phase structure.
\item[A7 -- universal symmetry:] supplies symmetry-governed relational laws and the primitive exchange action.
\item[A8 -- paradox stabilization:] lifts context-dependent incompatible projections into a higher-order closure carrier.
\end{description}

This partition creates a circularity firewall: minimal nontrivial distinction fixes binarity before exchange symmetry is introduced; A7 then acts on the already-binary relation. Arithmetic closure indices enter later through A5--A6.

\section{Minimal first distinction and primitive exchange orbit}

Let $\Omega_{\mathcal P}$ be the state/aspect domain carried by the point support. A distinction is a partition $\Pi$ of this domain. Nontriviality requires more than one nonempty block,
\begin{equation}
|\Pi|>1.
\end{equation}
Hence the least nontrivial distinction has exactly two outcomes,
\begin{equation}
\boxed{
|\Pi_1|=2,
\qquad
\Pi_1=\{N,S\}.
}
\label{eq:v12_minimal_binary_partition}
\end{equation}

\begin{theorem}[Minimal binary distinction]
Every nontrivial partition has at least two nonempty blocks. Therefore the minimal nontrivial distinction is binary.
\end{theorem}

The labels $N,S$ are relational poles or aspects of the minimal carrier. They are not assumed to be two pre-existing spatial objects.

Only after this binary distinction exists does the primitive exchange symmetry act:
\begin{equation}
\boxed{
J:N\leftrightarrow S,
\qquad
J^2=\mathrm{id}.
}
\label{eq:v12_primitive_involution}
\end{equation}
Choosing $x=N$ gives $Jx=S\neq x$ and
\begin{equation}
\boxed{
\mathcal O_J(x)=\{x,Jx\}=\{N,S\}.
}
\label{eq:v12_primitive_binary_orbit}
\end{equation}

Thus binarity is not inferred from an order-two involution. The involution is the exchange symmetry of the already-minimal binary distinction.
\vTwelveStatus{B}{--}{PASS}

\section{Unique exchange-invariant share}

Attach normalized nonnegative relational weights
\begin{equation}
w_N,w_S\ge0,
\qquad
w_N+w_S=1.
\end{equation}
Exchange invariance gives
\begin{equation}
(w_N,w_S)=(w_S,w_N),
\end{equation}
therefore
\begin{equation}
w_N=w_S.
\end{equation}
Normalization fixes the unique symmetric assignment,
\begin{equation}
\boxed{
(w_N,w_S)=\left(\frac12,\frac12\right).
}
\label{eq:v12_unique_half_share}
\end{equation}

\begin{theorem}[Unique exchange-invariant share]
The unique normalized nonnegative weight assignment invariant under exchange of the two primitive poles is $(1/2,1/2)$.
\end{theorem}

Introducing the scalar coordinate
\[
u=w_S,
\qquad
w_N=1-u,
\]
turns the pole exchange into
\begin{equation}
J(u)=1-u.
\end{equation}
Its unique fixed point is
\begin{equation}
\boxed{\operatorname{Fix}(J)=\left\{\frac12\right\}.}
\label{eq:v12_half_fixed_point}
\end{equation}
Thus the normalized exchange-invariant vector and the half coordinate are two descriptions of the same primitive balance.

\section{First information value}

For the normalized binary relation, use the Shannon carrier
\begin{equation}
H_2(u)=-(1-u)\ln(1-u)-u\ln u.
\end{equation}
At the exchange-fixed share,
\begin{equation}
\boxed{H_2(1/2)=\ln2.}
\label{eq:v12_first_ln2}
\end{equation}
This yields the first exact scalar dependency chain
\begin{equation}
\boxed{
\mathfrak Z_{\rm rel}
\rightarrow
\mathcal P
\rightarrow
\{N,S\}
\xrightarrow{\rm exchange\ invariance}
\frac12
\xrightarrow{H_2}
\ln2.
}
\label{eq:v12_primitive_scalar_chain}
\end{equation}
Its minimal TIR parent set is A0, A1, the minimal-nontrivial-distinction rule, A3 and A7. A2 enters only at the coherent quantum lift.

\section{Quantum lift}

Represent the two poles by an orthonormal basis
\[
|N\rangle,\qquad |S\rangle.
\]
Their minimal complex span is
\begin{equation}
\boxed{\mathcal H_{NS}\cong\mathbb C^2.}
\label{eq:v12_binary_c2}
\end{equation}
The balanced coherent family is
\begin{equation}
\boxed{
|\psi_{1/2}(\varphi)\rangle
=\frac{|N\rangle+e^{i\varphi}|S\rangle}{\sqrt2},
\qquad
\varphi\in\mathbb R/2\pi\mathbb Z.
}
\label{eq:v12_balanced_spinor}
\end{equation}
The relative phase is therefore an additional coherent coordinate above the same scalar half point.

The primitive output packet is
\begin{equation}
\boxed{
\mathfrak Z_{\rm rel}
\prec_{\rm dep}
\mathcal P
\prec_{\rm dep}
\{N,S\}
\prec_{\rm dep}
\frac12
\prec_{\rm dep}
\ln2
\prec_{\rm dep}
\mathbb C^2.
}
\label{eq:v12_primitive_packet}
\end{equation}

\section{Branch separation}

Projectivization of the two-state carrier and the spatial geometry developed in Parts I--II proceed from the common primitive packet. The three-flavour/Standard-Model branch proceeds from the same root after the geometric carrier is established. The sibling time programme receives the primitive packet through its own typed interface.

The dependency architecture is therefore
\[
\boxed{
\mathcal C_0
\prec_{\rm dep}
\left\{
\mathcal G_X^{\rm TIR},
\mathcal B_{\rm SM}^{\rm TIR},
\mathcal B_T^{\rm IDT}
\right\},
}
\]
where $\mathcal C_0$ denotes the admitted primitive information/quantum packet.

The spatial and temporal branches later meet through the declared closure surface
\begin{equation}
\mathfrak S_X^{\rm TIR}\otimes\mathfrak T^{\rm IDT}
\longrightarrow
\mathfrak M_{XT},
\end{equation}
and the matter branch then couples to the resulting spacetime carrier. The formal completion order of these joins is recorded in Chapter~\ref{ch:v12_completion_frontier}.

\section{First-chapter invariant}

The foundational narrative used by the rest of v12 is therefore
\begin{equation}
\boxed{
\text{relational zero}
\to\text{minimal point support}
\to\text{minimal nontrivial distinction}
\to\text{two poles}
\to\text{exchange symmetry}
\to\frac12
\to\ln2
\to\mathbb C^2.
}
\end{equation}
The next chapter resolves the geometry and information content of the half coordinate before any flavour normalization is introduced.

% ---- END TIR/monograph/v12/chapters/ch01_first_distinction_relational_kernel.tex ----
\clearpage

% ---- BEGIN TIR/monograph/v12/chapters/ch02_half_binary_information_ln2.tex ----
% !TEX root = ../../tir_monograph_v12.tex
\chapter[Half Coordinate, Binary Information and ln2]{The Half Coordinate, Binary Information and $\ln2$}
\label{ch:v12_half_ln2}

\section{Binary probability coordinate}

The primitive pair from Chapter~\ref{ch:v12_first_distinction} admits the normalized probability coordinate
\begin{equation}
\mathbf p(u)=(1-u,u),
\qquad
0\le u\le1.
\end{equation}
Pole exchange is
\begin{equation}
J(u)=1-u.
\end{equation}
The unique fixed point is
\begin{equation}
\boxed{u_\star=\frac12.}
\label{eq:v12_half_point}
\end{equation}
At this point
\begin{equation}
p_N=p_S=\frac12,
\qquad
\boxed{H_2=\ln2.}
\label{eq:v12_half_entropy}
\end{equation}
The scalar half and the first symmetric information value are therefore fixed before the flavour and particle branches enter.
\vTwelveStatus{B}{--}{PASS}

\section{Coherent lift of the half point}

A normalized pure state of the two-state carrier can be written, after removal of the global phase, as
\begin{equation}
|\psi(u,\varphi)\rangle
=\sqrt{1-u}\,|N\rangle
+e^{i\varphi}\sqrt{u}\,|S\rangle,
\qquad
\varphi\in\mathbb R/2\pi\mathbb Z.
\label{eq:v12_pure_binary_state}
\end{equation}
Projection to the scalar probability coordinate is
\begin{equation}
\pi:\ |\psi(u,\varphi)\rangle\longmapsto(1-u,u).
\end{equation}
Over the half point,
\begin{equation}
\boxed{
\pi^{-1}\!\left(\frac12,\frac12\right)
=
\left\{
\frac{|N\rangle+e^{i\varphi}|S\rangle}{\sqrt2}
:\varphi\in S^1
\right\}.
}
\label{eq:v12_half_fiber}
\end{equation}
Thus the scalar half coordinate lifts to a one-dimensional relative-phase fiber
\begin{equation}
\boxed{\frac12\longrightarrow U(1)\cong S^1.}
\end{equation}
This fiber carries coherent information that is absent from the probability projection.

\section{Bloch realization of the half fiber}

For Eq.~\eqref{eq:v12_pure_binary_state}, the Bloch vector is
\begin{equation}
\mathbf r(u,\varphi)
=\left(
2\sqrt{u(1-u)}\cos\varphi,
2\sqrt{u(1-u)}\sin\varphi,
1-2u
\right).
\end{equation}
Its norm satisfies
\begin{equation}
|\mathbf r|^2
=4u(1-u)+(1-2u)^2
=1.
\end{equation}
Hence pure states lie on $S^2$. At the half point,
\begin{equation}
\boxed{
\mathbf r\left(\frac12,\varphi\right)
=(\cos\varphi,\sin\varphi,0),
}
\label{eq:v12_half_equator}
\end{equation}
which is the Bloch equator.

The primitive two-pole relation therefore becomes
\begin{equation}
\boxed{
N\text{ pole}
\quad\longleftrightarrow\quad
S^1_{\rm half}\text{ equator}
\quad\longleftrightarrow\quad
S\text{ pole}.
}
\end{equation}
The projective geometry of this carrier is developed systematically in Chapter~\ref{ch:v12_complex_carrier}.

\section{Pole exchange on the phase fiber}

Let the exchange operator act by
\begin{equation}
X|N\rangle=|S\rangle,
\qquad
X|S\rangle=|N\rangle.
\end{equation}
On the balanced family,
\begin{equation}
X|\psi_{1/2}(\varphi)\rangle
\sim
|\psi_{1/2}(-\varphi)\rangle,
\end{equation}
where $\sim$ denotes projective equivalence after removal of a global phase. Hence
\begin{equation}
\boxed{\varphi\mapsto-\varphi\pmod{2\pi}.}
\label{eq:v12_half_phase_exchange}
\end{equation}
The exchange-fixed projective phases obey
\begin{equation}
\varphi=-\varphi\pmod{2\pi},
\end{equation}
so
\begin{equation}
\boxed{\varphi\in\{0,\pi\}.}
\label{eq:v12_half_fixed_phases}
\end{equation}
These are the two exchange eigen-directions on the equator.

\section{Phase closure and arithmetic indices}

For a closed phase loop,
\begin{equation}
\boxed{
n=\frac1{2\pi}\oint d\varphi\in\mathbb Z.
}
\label{eq:v12_winding}
\end{equation}
Equivalently, with $z=e^{i\varphi}$,
\begin{equation}
z^n=1.
\end{equation}
For a primitive $n$-fold closure,
\begin{equation}
\varphi=\frac{2\pi k}{n}.
\end{equation}
This supplies the operational A5--A6 chain
\begin{equation}
\boxed{
S^1
\longleftrightarrow
U(1)
\longrightarrow
\text{winding/closure invariant}
\longrightarrow
n\in\mathbb N.
}
\end{equation}
The ordering is important: the phase geometry is admitted before the arithmetic closure index is used as a TIR physical label.

\section{Information geometry at the half seam}

The binary probability projection sees the Bloch vertical coordinate
\begin{equation}
z_B=1-2u.
\end{equation}
At $u=1/2$,
\begin{equation}
z_B=0.
\end{equation}
All balanced relative phases therefore project to the same scalar probability point,
\begin{equation}
\boxed{
\pi:S^1_{\rm equator}\to\left\{\frac12\right\}.
}
\end{equation}
This is the exact interface between binary information and the coherent phase degree of freedom: the scalar coordinate fixes equal modulus, while the $U(1)$ fiber carries the relative phase.

\section{Dynamical interfaces}

A strongly continuous unitary flow on the half fiber admits a self-adjoint generator under the standard unitary-generator theorem. In the temporal branch, a parametrization
\begin{equation}
\varphi=\varphi(\tau)
\end{equation}
therefore supplies the entry point to phase dynamics. Multiple Bloch distinction axes similarly give noncommuting Pauli generators and the standard Robertson--Heisenberg structure.

These are downstream interfaces of the carrier established here. The present chapter fixes the static packet
\begin{equation}
\boxed{
\frac12
\leftrightarrow
\ln2
\leftrightarrow
S^1_{\rm phase}
\subset
\mathbb{CP}^1\cong S^2.
}
\end{equation}

\section{Ownership boundary}

Chapter 2 owns the binary half, $\ln2$, and its coherent phase fiber. Chapter 8 later introduces the three-flavour carrier, and Chapter 9 is the sole v12 publication owner of the complete normalization
\[
\kappa=\frac{\ln2}{24\pi}.
\]
This dependency order keeps the binary information numerator upstream of the flavour-mixing phase measure that normalizes it.

% ---- END TIR/monograph/v12/chapters/ch02_half_binary_information_ln2.tex ----
\clearpage

% ---- BEGIN TIR/monograph/v12/chapters/ch03_complex_carrier_information_geometry.tex ----
% !TEX root = ../../tir_monograph_v12.tex
\chapter{Complex Relational Carrier and Information Geometry}
\label{ch:v12_complex_carrier}

\section{Two-state carrier}

The primitive quantum lift gives
\begin{equation}
\mathcal H_2\cong\mathbb C^2.
\end{equation}
Choose the pole basis $|N\rangle,|S\rangle$. A normalized pure state can be parameterized by
\begin{equation}
\boxed{
|\psi(\vartheta,\varphi)\rangle
=\cos\frac\vartheta2\,|N\rangle
+e^{i\varphi}\sin\frac\vartheta2\,|S\rangle.
}
\label{eq:v12_spinor_spherical}
\end{equation}
The associated probability coordinate is
\begin{equation}
u=\sin^2\frac\vartheta2,
\qquad
1-u=\cos^2\frac\vartheta2.
\end{equation}
Hence the half seam $u=1/2$ is the equatorial condition
\begin{equation}
\vartheta=\frac\pi2.
\end{equation}

\section{Projective state space}

Physical pure-state rays remove global complex phase, giving
\begin{equation}
\boxed{\mathbb{CP}^1\cong S^2.}
\label{eq:v12_cp1_s2}
\end{equation}
The standard Fubini--Study line element for normalized rays may be written
\begin{equation}
 ds_{\rm FS}^2
=\langle d\psi|d\psi\rangle
-|\langle\psi|d\psi\rangle|^2
\label{eq:v12_fs_general}
\end{equation}
and, for Eq.~\eqref{eq:v12_spinor_spherical}, becomes
\begin{equation}
\boxed{
 ds_{\rm FS}^2
=\frac14\left(d\vartheta^2+\sin^2\vartheta\,d\varphi^2\right).
}
\label{eq:v12_fs_qubit}
\end{equation}
Thus the qubit projective metric is one quarter of the unit Bloch-sphere metric. Its total area is
\begin{equation}
\boxed{\operatorname{Area}_{FS}(\mathbb{CP}^1)=\pi.}
\end{equation}
This normalized projective area later provides an exact geometric measure for tetrahedral crosswalks; physical length/area calibration remains separately typed.

\section{Density-matrix representation}

For a normalized pure state,
\begin{equation}
\rho_\psi=|\psi\rangle\langle\psi|
\end{equation}
can be written in Bloch form
\begin{equation}
\boxed{
\rho_\psi
=\frac12\left(I+\mathbf r\cdot\boldsymbol\sigma\right),
\qquad
|\mathbf r|=1.
}
\label{eq:v12_bloch_density}
\end{equation}
The coefficient vector is
\begin{equation}
\mathbf r
=(\sin\vartheta\cos\varphi,
\sin\vartheta\sin\varphi,
\cos\vartheta).
\end{equation}
For a general normalized qubit state, $|\mathbf r|\le1$ and the trace-one Hermitian state space has affine hull
\begin{equation}
\boxed{
\mathcal A_2
=\{\rho=\rho^\dagger:\operatorname{Tr}\rho=1\}
=\frac12I+\operatorname{Herm}_0(2).
}
\label{eq:v12_affine_qubit_hull}
\end{equation}

\section{Information coordinate inside projective geometry}

The binary probability coordinate is related to the Bloch vertical coordinate by
\begin{equation}
r_z=1-2u=\cos\vartheta.
\end{equation}
The exchange-fixed information point therefore obeys
\begin{equation}
u=\frac12
\iff
r_z=0
\iff
\vartheta=\frac\pi2.
\end{equation}
The phase coordinate $\varphi$ remains free along the equator. This makes the information/projective relation explicit:
\begin{equation}
\boxed{
\text{binary share }u
\longleftrightarrow
\text{Bloch latitude},
\qquad
\text{relative phase }\varphi
\longleftrightarrow
\text{Bloch azimuth}.
}
\end{equation}

At the half seam the entropy is maximal,
\begin{equation}
H_2(1/2)=\ln2,
\end{equation}
while the coherent family still carries a full $S^1$ phase orbit. The scalar information coordinate and the projective phase coordinate are therefore complementary coordinates of the same two-state carrier.

\section{Operator coordinates}

The real vector space of Hermitian operators on $\mathbb C^2$ admits the unique Pauli decomposition
\begin{equation}
\boxed{
A=a_0I+a_x\sigma_x+a_y\sigma_y+a_z\sigma_z,
\qquad a_\mu\in\mathbb R.
}
\label{eq:v12_pauli_decomposition}
\end{equation}
with
\begin{equation}
\dim_{\mathbb R}\operatorname{Herm}(2)=4.
\end{equation}
The trace component is the uniform direction; the distinction-generating part lies in
\begin{equation}
\operatorname{Herm}_0(2)
=\operatorname{span}_{\mathbb R}\{\sigma_x,\sigma_y,\sigma_z\}.
\end{equation}
Chapter~\ref{ch:v12_c2_herm0_r3} turns this observation into the explicit three-real-dimensional relational generator carrier and its affine spatial relation object.

\section{Rotation covariance}

For $U\in SU(2)$,
\begin{equation}
A\mapsto UAU^\dagger
\end{equation}
preserves trace and the Hilbert--Schmidt pairing. On the Pauli coefficient vector the standard adjoint action induces
\begin{equation}
\boxed{
\operatorname{Ad}:SU(2)\to SO(3),
\qquad
SU(2)/\{\pm I\}\cong SO(3).
}
\label{eq:v12_su2_so3}
\end{equation}
The Bloch sphere, traceless operator sphere and rotation group are therefore coordinated by one algebra rather than introduced as unrelated copies of $S^2$ and $\mathbb R^3$.

\section{Fubini--Study tetrahedral preview}

Four unit Bloch vectors satisfying
\begin{equation}
\mathbf n_a\cdot\mathbf n_b=-\frac13
\qquad(a\neq b)
\end{equation}
form a regular tetrahedral frame. Their Bloch central angle obeys
\begin{equation}
\cos\chi=-\frac13.
\end{equation}
For the corresponding geodesic $\mathbb{CP}^1$ face, the spherical interior angle is $2\pi/3$ and the Fubini--Study area is
\begin{equation}
\boxed{a_{FS}^{\rm face}=\frac\pi4.}
\end{equation}
Four faces therefore cover total Fubini--Study area
\begin{equation}
\boxed{a_{FS}^{\rm tet}=\pi,}
\end{equation}
matching the full normalized $\mathbb{CP}^1$ area. Chapter~\ref{ch:v12_tetrahedral_closure} returns to this frame after the spatial carrier and Euclidean dimension gate have been established.

\section{Dependency output}

The chapter exports the exact mathematical packet
\begin{equation}
\boxed{
\mathbb C^2
\to
\mathbb{CP}^1\cong S^2
\to
\rho=\frac12(I+\mathbf r\cdot\boldsymbol\sigma)
\to
\operatorname{Herm}(2)
\supset
\operatorname{Herm}_0(2).
}
\end{equation}
The next chapter isolates the traceless translation/generator space and constructs the canonical affine relation
\[
\mathcal E_{xy}=2(\rho_y-\rho_x),
\]
which is the bridge from projective quantum-state geometry to the local three-real-dimensional relational carrier.

% ---- END TIR/monograph/v12/chapters/ch03_complex_carrier_information_geometry.tex ----
\clearpage
